\subsubsection{データ永続化}

\term{データ永続化}{Data Persistence} は、プログラム終了後もデータを保持し、必要なときに取り出せるようにすることである。設計では、データベース、ファイル、キャッシュ、メッセージキュー、外部ストレージなどをどう使うかを決める。

永続化の設計では、保存形式、整合性、トランザクション、検索性能、バックアップ、復旧、権限管理を考慮する。データは長期に残るため、後から変更しにくい。したがって、データ構造は実装の都合だけでなく、将来の変更と保守を見越して決める必要がある。
